\begin{tabularx}{\textwidth}{L{28mm}X X X}
\toprule
\textbf{Diagram} & \textbf{Fő kérdés} & \textbf{Fő elemek} & \textbf{Mikor hasznos?} \\
\midrule
Use case diagram & Mit tud a rendszer a külső szereplők számára? & Aktorok, használati esetek, rendszerhatár, include, extend, általánosítás. & Követelmények tisztázásánál, jogosultságok és felhasználói célok áttekintésénél. \\
Osztálydiagram & Milyen statikus fogalmakból, osztályokból és kapcsolatokból áll a rendszer? & Osztályok, attribútumok, operációk, asszociációk, öröklődés, aggregáció, kompozíció. & Elemzésnél, objektumorientált tervezésnél, domain modell és implementációs modell készítésénél. \\
Szekvenciadiagram & Egy konkrét forgatókönyvben milyen sorrendben kommunikálnak az objektumok? & Résztvevők, életvonalak, aktivációk, üzenetek, válaszok. & Használati esetek részletezésénél, felelősségek és időrend tisztázásánál. \\
Komponensdiagram & Milyen nagyobb szoftverkomponensekből áll a rendszer, és hogyan függenek egymástól? & Komponensek, interfészek, függőségek, portok. & Architektúra, implementációs egységek, integrációs pontok és telepítési előkészítés áttekintésénél. \\
\bottomrule
\end{tabularx}
